% Архитектура Drone Swarm Simulator v2.
% Сборка: pdflatex drone_swarm_simulator_architecture.ru.tex
%
% Основные источники:
%   Drone_Swarm_Simulator_v2/README.md
%   Drone_Swarm_Simulator_v2/launch_simulation.py
%   Drone_Swarm_Simulator_v2/core/control/drone_controller.py
%   Drone_Swarm_Simulator_v2/core/mavlink/worker.py
%   Drone_Swarm_Simulator_v2/scenarios/*.py
%   Drone_Swarm_Simulator_v2/visualizer/*.py

\documentclass[a4paper,11pt]{article}

\usepackage[T2A]{fontenc}
\usepackage[utf8]{inputenc}
\usepackage[russian]{babel}
\usepackage{amsmath,amssymb}
\usepackage{geometry}
\usepackage{hyperref}
\usepackage{enumitem}
\usepackage{array}
\geometry{left=2.5cm,right=2.5cm,top=2.5cm,bottom=2.5cm}
\setlist{nosep,leftmargin=*}
\hypersetup{colorlinks=true,linkcolor=blue,urlcolor=blue,citecolor=blue}

\newcommand{\code}[1]{\texttt{\detokenize{#1}}}
\newcommand{\MAVLink}{\textsc{MAVLink}}
\newcommand{\SITL}{\textsc{SITL}}

\title{Архитектура симулятора роя БПЛА\\
\code{Drone\_Swarm\_Simulator\_v2}}
\author{}
\date{}

\begin{document}
\maketitle

\begin{abstract}
Документ описывает устройство симулятора \code{Drone\_Swarm\_Simulator\_v2}:
его назначение, основные программные блоки, модель обмена данными между ними
и роль пользовательского сценария в управлении роем. Симулятор построен как
набор взаимодействующих процессов вокруг ArduPilot \SITL: каждый виртуальный
дрон исполняется отдельным экземпляром автопилота, а пользовательский
Python-сценарий подключается к каждому дрону по протоколу \MAVLink, читает
телеметрию, рассчитывает закон управления и отправляет команды обратно.
\end{abstract}

\tableofcontents
\newpage

%====================================================================
\section{Назначение симулятора}
%====================================================================

\code{Drone\_Swarm\_Simulator\_v2} предназначен для исследования алгоритмов
управления группой однородных беспилотных летательных аппаратов без
необходимости запускать реальные дроны. В отличие от чисто математической
модели, здесь в контуре присутствует полноценный автопилот ArduPilot:
команды управления проходят через те же интерфейсы, что и при работе с
реальным аппаратом, а состояние аппарата получается из телеметрии
автопилота.

Основные цели симулятора:
\begin{enumerate}
  \item \textbf{Проверка ройевых алгоритмов.} Сценарии позволяют задавать
  поведение лидеров, ведомых, цепочек, ``антенн'' и других формаций, а затем
  наблюдать, сохраняет ли рой требуемую геометрию.

  \item \textbf{Отработка взаимодействия с автопилотом.} Управление идёт не
  напрямую через координаты физической модели, а через команды \MAVLink:
  армирование, взлёт, смена режима, \code{RC\_CHANNELS\_OVERRIDE} и запросы
  телеметрии. Это приближает эксперимент к реальному стеку управления.

  \item \textbf{Сбор воспроизводимых данных.} Запуски сохраняются в CSV и
  \code{metadata.json}; эти данные можно анализировать графиками или
  воспроизводить в RViz.

  \item \textbf{Сравнение идеального и оценочного состояния.} Часть сценариев
  использует одновременно истинную позу симулятора и оценку автопилота
  (\code{SIM\_STATE}, \code{LOCAL\_POSITION\_NED},
  \code{GLOBAL\_POSITION\_INT}), что позволяет исследовать влияние шума и
  навигационной оценки на закон управления.
\end{enumerate}

По смыслу данный симулятор является исследовательской платформой для
итеративной разработки: сначала алгоритм можно быстро проверить в
SITL-only-режиме, затем при необходимости подключить 3D-визуализацию или
воспроизвести результат в ROS/RViz для анализа траекторий.

%====================================================================
\section{Общая архитектура}
%====================================================================

Симулятор состоит из четырёх основных блоков:
\begin{enumerate}
  \item \textbf{Лаунчер} \code{launch\_simulation.py} --- главный управляющий
  скрипт, который запускает остальные процессы.

  \item \textbf{ArduPilot \SITL} --- набор экземпляров автопилота, по одному
  процессу на каждый виртуальный дрон.

  \item \textbf{Сценарий} из каталога \code{scenarios/} --- пользовательский
  алгоритм, который соединяется с дронами, обменивается состояниями внутри
  роя, вычисляет управляющие воздействия и пишет логи.

  \item \textbf{Визуализация и воспроизведение} --- 2D-окно matplotlib во
  время эксперимента, опциональный Webots-режим и офлайн-воспроизведение в
  RViz по CSV-логам.
\end{enumerate}

Упрощённая схема запуска и обмена данными:

\begin{verbatim}
launch_simulation.py
    |
    | subprocess
    v
ArduPilot SITL #1  <--- MAVLink UDP --->  MAVLinkWorker #1
ArduPilot SITL #2  <--- MAVLink UDP --->  MAVLinkWorker #2
...                                      ...
ArduPilot SITL #N  <--- MAVLink UDP --->  MAVLinkWorker #N
                                           |
                                           v
                                  сценарий scenarios/*.py
                                           |
                 +-------------------------+-------------------------+
                 |                         |                         |
                 v                         v                         v
          CSV/metadata logs          UDP JSON 2D viz          swarm logic
\end{verbatim}

Важно, что рой как распределённая система моделируется на уровне сценария.
Каждый дрон имеет собственный экземпляр ArduPilot \SITL, но обмен позициями
между дронами в большинстве сценариев выполняется не через радиоканал и не
через междронный \MAVLink, а через структуры данных внутри одного
Python-процесса сценария. Это сделано для удобства исследования законов
управления: физика и автопилот остаются раздельными для каждого аппарата, а
алгоритм роя получает доступ к актуальным координатам соседей через общий
обменный цикл.

%====================================================================
\section{Лаунчер как верхний уровень управления}
%====================================================================

Файл \code{Drone\_Swarm\_Simulator\_v2/launch\_simulation.py} является единой
точкой входа. Он не реализует сам закон управления, а выполняет роль
оркестратора:
\begin{enumerate}
  \item разбирает аргументы командной строки: режим запуска, количество
  дронов, имя сценария, длительность, директорию эксперимента, частоту обмена;
  \item при необходимости запускает Webots;
  \item запускает требуемое число экземпляров \code{sim\_vehicle.py};
  \item ждёт \code{HEARTBEAT} от каждого ArduPilot \SITL, чтобы убедиться, что
  MAVLink-порты готовы;
  \item запускает 2D-визуализатор, если он выбран;
  \item запускает сценарий как отдельный Python-процесс;
  \item при завершении сценария или при сигнале остановки завершает дочерние
  процессы.
\end{enumerate}

Сценарии регистрируются в списке \code{SCENARIOS}. Для каждого сценария
задаются идентификатор, человекочитаемое описание, путь к файлу и рабочая
директория. Отдельный словарь \code{SCENARIO\_MIN\_DRONES} задаёт минимальное
количество аппаратов для сценариев, которым требуется больше одного дрона
например, для линейной цепочки или антенны.

Типичный запуск выглядит так:
\begin{verbatim}
source ../drone_env/bin/activate
python launch_simulation.py -s -c fntena_logic_copy2 -n 4 --duration 60
\end{verbatim}

Здесь флаг \code{-s} выбирает режим SITL-only, \code{-c} задаёт сценарий,
\code{-n} --- число дронов, а \code{--duration} --- длительность эксперимента.

%====================================================================
\section{Блок ArduPilot \SITL}
%====================================================================

ArduPilot \SITL~--- это программная версия автопилота. Она исполняет тот же
код управления, что и реальный ArduPilot, но вместо физических датчиков и
моторов использует симулированное окружение. В данном проекте каждый дрон
запускается отдельным процессом \code{sim\_vehicle.py}.

Для каждого аппарата лаунчер задаёт:
\begin{itemize}
  \item \code{--instance}: номер экземпляра \SITL;
  \item \code{--sysid}: MAVLink-идентификатор дрона;
  \item домашнюю точку \code{--custom-location}: широта, долгота, высота и
  курс;
  \item файл параметров \code{config/iris.parm};
  \item набор выходных MAVLink-портов через \code{--out}.
\end{itemize}

Домашние точки дронов немного смещаются по восточной координате. Благодаря
этому у всех аппаратов получается общий локальный NED-кадр, где дроны
изначально стоят не в одной точке, а с небольшим шагом по оси East.

Основной канал связи между сценарием и \SITL --- \MAVLink поверх UDP. Порты
распределяются по простой схеме:

\begin{center}
\begin{tabular}{>{\raggedright\arraybackslash}p{0.38\textwidth}
                >{\raggedright\arraybackslash}p{0.30\textwidth}
                >{\raggedright\arraybackslash}p{0.18\textwidth}}
\hline
Назначение & Формула порта для дрона $i$ & Пример для $i=0$ \\
\hline
Основной канал сценария & $14551 + 10i$ & 14551 \\
Пассивный канал логгера & $14651 + 10i$ & 14651 \\
Дополнительный raw/tlog-канал & $14751 + 10i$ & 14751 \\
\hline
\end{tabular}
\end{center}

Такое разделение нужно, чтобы основной сценарий и пассивные инструменты
логирования не конкурировали за один и тот же UDP-поток.

%====================================================================
\section{MAVLink-слой и контроллер дрона}
%====================================================================

В сценарии прямой доступ к \code{pymavlink} не размазывается по всему коду.
Для каждого аппарата создаётся объект \code{DroneController}, а внутри него
--- \code{MAVLinkWorker}. Это два разных уровня абстракции.

\subsection{MAVLinkWorker}

\code{MAVLinkWorker} отвечает за потокобезопасный обмен с одним дроном.
Все операции \code{recv\_msg()} и отправки MAVLink-команд выполняются в одном
выделенном daemon-потоке. Остальная часть сценария не читает UDP-сокет
напрямую, а кладёт команды в очередь и читает уже обновлённый кэш состояния.

Работник выполняет несколько задач:
\begin{itemize}
  \item подключается к строке вида \code{udp:127.0.0.1:14551};
  \item ждёт \code{HEARTBEAT};
  \item запрашивает частоту сообщений \code{HOME\_POSITION},
  \code{SIM\_STATE} и \code{LOCAL\_POSITION\_NED};
  \item принимает телеметрию и обновляет последнюю позицию, ориентацию и
  оценочную позицию;
  \item исполняет команды из очереди: смена режима, армирование, взлёт,
  \code{RC\_CHANNELS\_OVERRIDE}.
\end{itemize}

Позиция для основного контура обычно берётся из \code{SIM\_STATE}: это
``истинное'' состояние симулятора. Для задач сравнения может использоваться
оценка автопилота из \code{LOCAL\_POSITION\_NED} или
\code{GLOBAL\_POSITION\_INT}.

\subsection{DroneController}

\code{DroneController} --- более удобная обёртка для сценариев. Он хранит
конфигурацию дрона, создаёт \code{MAVLinkWorker}, предоставляет методы
\code{get\_position()}, \code{initialize()}, \code{send\_rc\_override()} и
запускает поток keepalive, который периодически отправляет последние
RC-команды. Это важно, потому что режим управления через RC override требует
регулярного обновления каналов.

Кроме собственного состояния дрона, \code{DroneController} хранит словарь
\code{other\_drones\_positions}. Именно через него сценарий передаёт одному
виртуальному дрону координаты остальных участников роя. С точки зрения
архитектуры это не отдельный сетевой протокол, а внутрипроцессная модель
обмена: сценарий читает позиции всех дронов и раскладывает их по словарям
контроллеров.

%====================================================================
\section{Скрипт сценария}
%====================================================================

Скрипт сценария --- это место, где задаётся смысл эксперимента. Лаунчер только
поднимает инфраструктуру, ArduPilot исполняет низкоуровневую логику полёта, а
сценарий определяет, что должен делать рой.

Типичный сценарий выполняет следующие шаги:
\begin{enumerate}
  \item читает аргументы командной строки, которые передал лаунчер;
  \item формирует список конфигураций дронов:
  \code{\{"id": 1, "udp\_port": 14551, "role": "..."\}};
  \item создаёт \code{DroneController} для каждого дрона;
  \item параллельно подключается к MAVLink-портам и выполняет начальную
  последовательность: армирование, взлёт, переход в нужный режим;
  \item запускает обменный цикл координат;
  \item запускает управляющие циклы для отдельных ролей дронов;
  \item пишет CSV-логи и метаданные эксперимента;
  \item по завершении переводит аппараты в безопасное состояние и закрывает
  MAVLink-работников.
\end{enumerate}

В сценариях роя обычно есть два типа циклов:
\begin{itemize}
  \item \textbf{exchange loop} --- общий цикл обмена. Он регулярно читает
  позиции всех контроллеров, обновляет \code{other\_drones\_positions},
  публикует координаты в 2D-визуализатор и пишет строки в CSV;

  \item \textbf{control loop} --- управляющий цикл конкретного дрона или
  группы дронов. Он читает собственную позицию и позиции соседей, вычисляет
  ошибку формации, значение $\sigma$, PID-регулятор или другой закон, после
  чего отправляет MAVLink-команду управления.
\end{itemize}

Именно сценарий задаёт модель ``общения'' роя на алгоритмическом уровне. Если
нужно смоделировать идеальный обмен координатами, сценарий напрямую передаёт
координаты соседей. Если нужно исследовать задержку, шум или ограниченную
видимость, это также реализуется в сценарии: например, можно фильтровать
соседей по радиусу, добавлять шум к координатам или брать не истинное
состояние \code{SIM\_STATE}, а оценочную навигационную позицию.

%====================================================================
\section{Визуализатор и воспроизведение}
%====================================================================

В проекте есть несколько уровней визуализации, и они решают разные задачи.

\subsection{2D-визуализатор matplotlib}

Оперативная визуализация во время запуска реализована в каталоге
\code{visualizer/}. Сценарий вызывает функцию \code{publish\_positions()},
которая отправляет JSON-пакет по UDP на порт \code{15551}. Отдельный процесс
\code{drone\_position\_visualizer.py} слушает этот порт, принимает текущие
координаты дронов в NED-кадре и рисует их на плоскости.

Формат связи простой:
\begin{verbatim}
{
  "positions": {
    "1": {"x": 0.0, "y": 0.0, "z": -1.0},
    "2": {"x": 1.0, "y": 0.0, "z": -1.0}
  },
  "rates": {
    "follower_hz": 50.0,
    "exchange_hz": 50.0,
    "webots_step_hz": null
  }
}
\end{verbatim}

Это односторонняя связь: визуализатор только принимает состояние и ничего не
отправляет обратно в сценарий. Если визуализатор не запущен, UDP-пакеты просто
теряются, а симуляция продолжается.

\subsection{Webots-режим}

Опциональный режим Webots используется как 3D-среда живой симуляции. В этом
режиме лаунчер сначала запускает Webots, затем запускает ArduPilot \SITL с
моделью \code{webots-python}. Обмен между Webots и \SITL идёт по TCP-портам
вида \code{5770 + 10i}. Webots отвечает за 3D-сцену и физическое окружение, а
ArduPilot продолжает быть автопилотом каждого аппарата.

В основном SITL-only-режиме Webots не нужен: ArduPilot сам предоставляет
симулированное состояние, достаточное для отладки MAVLink-управления и
ройевых законов.

\subsection{RViz replay}

RViz используется не как обязательная часть живой симуляции, а как средство
офлайн-воспроизведения. Сценарий сохраняет CSV-файлы и \code{metadata.json},
после чего скрипты из каталога \code{replay/} публикуют эти данные в ROS/ROS2:
позы дронов, маркеры, подписи и служебную информацию. Такой режим удобен для
анализа уже выполненного эксперимента и подготовки иллюстраций.

Если говорить строго о \code{Drone\_Swarm\_Simulator\_v2}, Gazebo не является
основным компонентом его рантайма. В проекте используются ArduPilot \SITL,
опциональный Webots, 2D-визуализатор matplotlib и офлайн-воспроизведение в
RViz.

%====================================================================
\section{Модель связи между блоками}
%====================================================================

Ниже сведены основные каналы обмена данными:

\begin{center}
\begin{tabular}{>{\raggedright\arraybackslash}p{0.30\textwidth}
                >{\raggedright\arraybackslash}p{0.30\textwidth}
                >{\raggedright\arraybackslash}p{0.28\textwidth}}
\hline
Связь & Тип связи & Назначение \\
\hline
Лаунчер $\rightarrow$ ArduPilot \SITL & \code{subprocess.Popen} &
Запуск процессов \code{sim\_vehicle.py}, передача CLI-параметров \\
\hline
Лаунчер $\rightarrow$ сценарий & \code{subprocess.Popen} &
Запуск выбранного Python-сценария с параметрами эксперимента \\
\hline
Сценарий $\leftrightarrow$ ArduPilot \SITL & \MAVLink поверх UDP &
Телеметрия, режимы, армирование, взлёт, RC override \\
\hline
Webots $\leftrightarrow$ ArduPilot \SITL & \MAVLink/TCP &
3D-физика и состояние аппарата в Webots-режиме \\
\hline
Сценарий $\rightarrow$ 2D-визуализатор & UDP + JSON &
Отображение текущих координат и частот циклов \\
\hline
Сценарий $\rightarrow$ CSV/metadata & Файловая запись &
Логи эксперимента для анализа и replay \\
\hline
Replay $\rightarrow$ RViz & ROS/ROS2 topics &
Офлайн-воспроизведение траекторий \\
\hline
Дрон $\leftrightarrow$ дрон внутри сценария & Python-словари в памяти &
Передача позиций соседей в алгоритм роя \\
\hline
\end{tabular}
\end{center}

Главный рабочий протокол управления --- \MAVLink. Именно через него сценарий
видит каждый виртуальный аппарат как отдельный автопилот. UDP JSON для
визуализации, файловые CSV-логи и ROS-топики для replay являются
вспомогательными каналами: они не управляют дроном, а помогают наблюдать и
анализировать эксперимент.

%====================================================================
\section{Роли блоков в одном эксперименте}
%====================================================================

В одном запуске обязанности распределены следующим образом:

\begin{description}
  \item[ArduPilot \SITL] моделирует автопилот конкретного БПЛА, принимает
  MAVLink-команды и выдаёт телеметрию. Это нижний уровень, близкий к реальному
  полётному стеку.

  \item[\code{MAVLinkWorker}] является техническим адаптером между Python-кодом
  сценария и MAVLink-соединением. Он изолирует работу с сокетом и делает
  доступ к состоянию потокобезопасным.

  \item[\code{DroneController}] представляет один дрон на уровне сценария:
  позволяет читать позицию, отправлять RC-команды, хранить позиции соседей и
  запускать начальную последовательность.

  \item[Сценарий] реализует исследуемую логику: роли дронов, закон формации,
  параметры соседства, частоту обмена, логирование и условия завершения.

  \item[Визуализатор] показывает состояние эксперимента человеку. Он не
  является частью контура управления и может быть отключён без изменения
  результата симуляции.

  \item[Лаунчер] соединяет все части в один запуск: создаёт процессы,
  проверяет готовность SITL и корректно завершает дочерние процессы.
\end{description}

Таким образом, архитектура разделяет ``что делать'' и ``как лететь''.
Сценарий отвечает за высокоуровневую ройевую цель, ArduPilot --- за исполнение
команд на уровне автопилота, а MAVLink-слой является интерфейсом между ними.

%====================================================================
\section{Задачи, которые решает такая архитектура}
%====================================================================

Выбранная архитектура полезна для исследовательского симулятора по нескольким
причинам.

\begin{enumerate}
  \item \textbf{Модульность.} Новый сценарий можно добавить как отдельный
  файл в \code{scenarios/}, не переписывая лаунчер, MAVLink-слой и
  визуализатор.

  \item \textbf{Масштабируемость по числу дронов.} Каждый дрон имеет свой
  процесс \SITL и свой UDP-порт. Добавление аппаратов сводится к запуску
  дополнительных экземпляров и созданию дополнительных контроллеров.

  \item \textbf{Повторяемость экспериментов.} Параметры запуска передаются
  через CLI, а результаты сохраняются в структурированном виде. Это позволяет
  запускать серии экспериментов и сравнивать результаты.

  \item \textbf{Близость к реальному управлению.} Алгоритм работает не с
  абстрактной точкой напрямую, а через MAVLink-команды к автопилоту. Поэтому
  проявляются задержки, частоты телеметрии, режимы автопилота и ограничения
  RC-управления.

  \item \textbf{Разделение исследования и отображения.} Визуализатор и RViz
  можно включать для анализа, но они не влияют на управляющий контур. Это
  уменьшает риск, что графический интерфейс изменит результат эксперимента.
\end{enumerate}

В итоге симулятор нужен не просто для ``красивого полёта'' нескольких моделей,
а для проверки гипотез о распределённом управлении: насколько устойчиво рой
держит формацию, как влияет качество навигации, как меняется поведение при
разных частотах обмена, какие ограничения появляются при переходе от
математической модели к автопилоту ArduPilot.

%====================================================================
\section{Краткий вывод}
%====================================================================

\code{Drone\_Swarm\_Simulator\_v2} можно рассматривать как связку из
оркестратора, набора автопилотов ArduPilot \SITL, пользовательского
Python-сценария и инструментов визуализации. Основной контур управления
выглядит так: сценарий получает телеметрию каждого дрона по \MAVLink,
обновляет внутреннюю модель соседей, вычисляет управляющее воздействие и
отправляет команды обратно в соответствующий экземпляр \SITL. Визуализация и
логирование существуют рядом с этим контуром и помогают анализировать
результат, но не подменяют саму модель управления.

\end{document}
